\begin{tikzpicture}[x=5.2cm,y=1.05cm,lab/.style={font=\tiny,align=center},ar/.style={-{Latex[length=1.4mm]},thick}]
\fill[black!4] (0,0) rectangle (1,2.18);
\fill[black!8] (0.28,0) rectangle (0.39,2.18);
\fill[black!8] (0.445,0) rectangle (0.555,2.18);
\fill[black!8] (0.61,0) rectangle (0.72,2.18);
\draw[black!45,thick] (0,0.50)--(1,0);
\draw[black!45,thick,densely dashed] (0,0)--(1,1);
\draw[black!45,densely dotted] (0.3333,0)--(0.3333,0.3333);
\fill[black!45] (0.3333,0.3333) circle (0.8pt);
\draw[black!65,thick] (0,1.00)--(1,0);
\draw[black!65,thick,densely dashed] (0,0)--(1,1);
\draw[black!65,densely dotted] (0.5000,0)--(0.5000,0.5000);
\fill[black!65] (0.5000,0.5000) circle (0.8pt);
\draw[black,very thick] (0,2.00)--(1,0);
\draw[black,very thick,densely dashed] (0,0)--(1,1);
\draw[black,densely dotted] (0.6667,0)--(0.6667,0.6667);
\fill[black] (0.6667,0.6667) circle (1.0pt);
\draw[-{Latex[length=1.4mm]}] (0,0)--(0,2.32) node[above,font=\scriptsize] {flux [arb.]};
\draw[-{Latex[length=1.4mm]}] (0,0)--(1.10,0) node[below,font=\scriptsize] {$\xi$};
\draw (0.3333,0.035)--(0.3333,-0.035) node[below,font=\tiny] {$1/3$};
\draw (0.5000,0.035)--(0.5000,-0.035) node[below,font=\tiny] {$1/2$};
\draw (0.6667,0.035)--(0.6667,-0.035) node[below,font=\tiny] {$2/3$};
\draw (1.0000,0.035)--(1.0000,-0.035) node[below,font=\tiny] {$1$};
\node[lab,anchor=west] at (0.04,2.03) {forward $r^+(1-\xi)$};
\node[lab,anchor=west] at (0.70,0.86) {reverse $r^-\xi$};
\node[lab] at (0.20,0.70) {$A<0$};
\node[lab] at (0.50,0.88) {$A=0$};
\node[lab] at (0.78,1.24) {$A>0$};
\draw[ar] (0.44,1.55)--(0.63,0.82) node[midway,right,font=\tiny,align=left] {$\xi_\eq$ moves\\with affinity};
\end{tikzpicture}
\caption{정지점 유도 개선안. 정방향 플럭스 $r^+(1-\xi)$ 와 역방향 플럭스 $r^-\xi$ 의 교점이 $\xi_\eq=r^+/(r^++r^-)$ 를 정한다(식~\eqref{eq:logisticsolve}). $r^+/r^-=2$ 예시는 $\xi_\eq=2/3$, detailed balance 기준 $r^+/r^-=1$ 은 $1/2$, 역구동 예시는 $1/3$ 으로 이동한다. 좌표는 세 비율의 직선과 교점을 Python 으로 수치 평가한 것이다.}
\label{fig:flux}
